\section{Problem Setting and Available Information}\label{sec:problem}

\subsection{Ordered policy profiles}

Let $h:[0,1]\to[0,1]$ denote an ordered policy profile.  A task $q$ supplies
ordered nodes
\[
  0<t_{q,1}<\cdots<t_{q,d_q}<1
\]
and an integer implementation scale $L_q$.  The deployed vector is
\[
  x_{q,i}(h)=\operatorname{round}\{L_qh(t_{q,i})\},
  \qquad i=1,\ldots,d_q.
\]
Thus $d_q$ is a grid resolution, not by itself the intrinsic dimension of the
admissible profile class.  Regular grids use
$t_{q,i}=(i-1/2)/d_q$; irregular grids are allowed when their ordering is
declared.  A schema may additionally declare channel order.  No claim is made
for an unordered vector with unknown coordinate semantics.

One simulator replication at $x$ returns an objective observation $Y^f_q(x)$
and a scalar constraint observation $Y^g_q(x)$.  We minimize
\[
  f_q(x)=\E[Y^f_q(x)]
\]
subject to the chance constraint
\begin{equation}\label{eq:chance-constraint}
  p_q(x):=\Prob\{Y^g_q(x)\le \tau_q\}\ge 1-\alpha_q.
\end{equation}
The true feasible set is $\cF_q=\{x:p_q(x)\ge1-\alpha_q\}$.
In synthetic experiments, the experimenter reveals truth only after design,
search, shortlisting, and verification are complete, and uses it only to score
the finished decision.

\subsection{Source and target information}

Let $q_1,\ldots,q_m$ be source tasks and $q_\star$ a target task.  Before the
target opens, every source task evaluates the same finite normalized profile
library $\cL=\{h_1,\ldots,h_J\}$.  For each pair $(q,h)$, the source archive
contains $R$ ordinary replications of $(Y^f_q,Y^g_q)$.  The primary method uses
$m=2$, $J=64$, and $R=3$, hence
\begin{equation}\label{eq:source-cost}
  S=mJR=384
\end{equation}
source simulator calls.  The library is constructed without source or target
outcomes; source data only determine which members are selected.

The source archive is supplied before the target problem opens.  In the
synthetic stress test its tasks come from the same declared regime as the
target; in the external study they obey the stated region-holdout rule.  The
method ranks profiles within this archive.  Choosing or purchasing related
historical tasks from a larger repository is a separate upstream decision.

Before target search, the method may use the target's ordered grid, integer
bounds, nominal dimension, and, in declared-schema experiments, channel order.
It may not use a target outcome, feasibility label, noise parameter, optimum,
safe-basin center, or verification sample.  The primary source aggregation is
domain blind: source tasks receive equal weight.  Descriptor-conditioned
weighting is evaluated only as a control.

\subsection{Decision stages and costs}

The initial-design step selects $n_0$ target policies.  A target optimizer may
then use a total of $N\ge n_0$ search calls and return an ordered list of at
most $K_s$ distinct candidates.  This list and its order are fixed before
fresh verification simulations begin.  The final outcome is either a
certified deployment or abstention.

We report the costs
\begin{equation}\label{eq:cost-identity}
  C_{\rm all}=S+N+V,
  \qquad
  C_{\rm amort}(M)=\frac{S}{M}+N+V,
\end{equation}
where $V$ is the realized independent-verification cost and $M$ is the number
of target deployments sharing one source archive.  ``Equal preverification
cost'' means that $S+N$ is matched before $V$; it is not called equal total
cost.

\subsection{Endpoints}

The primary search endpoint is whether the $N$ target evaluations contain a
truly feasible policy.  The primary deployment endpoint is whether the
verifier returns a truly feasible policy.  We also report false certificates,
verification calls, and a penalized loss that retains failures in the
denominator.  Conditional regret is reported only when a feasible policy
exists and is not used alone to compare methods with different feasible
counts.  A separate fixed library defines synthetic regret and an unattainable
finite-library upper bound; neither affects an implementable decision.
